\centerline{\bf \Huge ДЗ. Сумма цифр числа}
\centerline{\bf \Huge }

\begin{flushright}
        \scriptsize{\textbf{qq}}
\end{flushright}

\textbf{\textit{$S(X)$ - сумма цифр натурального числа $X$}}

\problem{\textbf{1}} Решите уравнение:

$X + S(X) + S(S(X)) = 100000000000000000000000000$

\problem{\textbf{2}} Решите уравнение:

$X + 2\cdot S(X) + 3\cdot S(S(X)) + 4 \cdot S(S(S(X)))  = 7777777$

\problem{\textbf{3}} Начальник написал на доске натуральные числа $N$ и $N+1$. Он заметил, что сумма цифр каждого из этих чисел делится на 13. Оказалось, что $N$ - наименьшее число, для которого это условие выполняется. Найдите $S(N)$, $S(N+1)$ и сами числа $N$, $N+1$.


\problem{\textbf{4}} Начальник написал на доске натуральные числа $N$ и $N+7$. Он заметил, что сумма цифр каждого из этих чисел делится на 101. Оказалось, что $N$ - наименьшее число, для которого это условие выполняется. Найдите $S(N)$, $S(N+7)$ и сами числа $N$, $N+7$.
